\documentclass[12pt,a4paper]{article}
\usepackage{estilo/icmc}
\addbibresource{../referencias.bib}
\graphicspath{{../dissertacao/figuras/}}  % figuras geradas por analise/viabilidade/ e analise/real/
% numeros da analise descritiva (gerados por analise/real/descritiva.py e r/descritiva.R)
% Gerado por analise/real/descritiva.py: nao editar a mao.
\newcommand{\descZeros}{\tr{88}{88}}
\newcommand{\descVarMedia}{\tr{1,40}{1.40}}
\newcommand{\descGiniObs}{\tr{0,72}{0.72}}
\newcommand{\descGiniPoisson}{\tr{0,69}{0.69}}
\newcommand{\descGiniHoras}{\tr{0,47}{0.47}}
\newcommand{\descGiniSiteObs}{\tr{0,68}{0.68}}
\newcommand{\descGiniSitePoisson}{\tr{0,65}{0.65}}
\newcommand{\descForaIC}{\tr{8}{8}}
\newcommand{\descPnoventaPdez}{\tr{85}{85}}
\newcommand{\descVelSpotter}{\tr{0,12}{0.12}}
\newcommand{\descVelShuttle}{\tr{0,48}{0.48}}
\newcommand{\descTurnoUm}{\tr{38}{38}}
\newcommand{\descTurnoDois}{\tr{36}{36}}
\newcommand{\descTurnoTres}{\tr{26}{26}}
\newcommand{\descTopCausas}{\tr{41}{41}}
\newcommand{\descSuspeitos}{\tr{24}{24}}
\newcommand{\descComHora}{\tr{97}{97}}
\newcommand{\descHorasExtras}{\tr{10}{10}}
\newcommand{\descNotaAud}{\tr{94}{94}}
\newcommand{\descAcidentesPct}{\tr{87}{87}}

% Gerado por r/descritiva.R: nao editar a mao.
\newcommand{\descRzerosObs}{\tr{88}{88}}
\newcommand{\descRzerosPoisson}{\tr{87}{87}}
\newcommand{\descRzerosNB}{\tr{88}{88}}
\newcommand{\descRzerosRE}{\tr{88}{88}}
\newcommand{\descRveicZeroObs}{\tr{53}{53}}
\newcommand{\descRveicZeroPoisson}{\tr{50}{50}}
\newcommand{\descRveicZeroRE}{\tr{54}{54}}
\newcommand{\descRdpVeiculo}{\tr{0,51}{0.51}}
\newcommand{\descRcorMax}{\tr{0,47}{0.47}}
\newcommand{\descRcorAcidMax}{\tr{0,11}{0.11}}
\newcommand{\descRcorHorasMax}{\tr{0,14}{0.14}}


% Compilar dentro desta pasta. O mesmo codigo gera as duas linguas:
%   portugues:  TEXINPUTS=".;../;" latexmk -pdf -jobname=PreProjeto_MECAI_Felipe_Terrazas main.tex
%   ingles:     TEXINPUTS=".;../;" latexmk -pdf -jobname=PreProjeto_MECAI_Felipe_Terrazas_EN main.tex
% Acrescente _PB no fim (..._PB ou ..._EN_PB) para a versao preto e branco.
% Convencao de escrita: ver README.md desta pasta.

\newcommand{\titulo}{\tr%
  {Risco de acidentes em operações de pátio: modelos hierárquicos de contagem com exposição e telemetria veicular}%
  {Accident risk in yard operations: hierarchical count models with exposure and vehicle telematics}}

\hypersetup{pdftitle={\titulo}}

\begin{document}

\begin{titlepage}
\centering
{\large\bfseries \tr{UNIVERSIDADE DE SÃO PAULO}{UNIVERSITY OF SÃO PAULO}\par}
{\large\bfseries \tr{INSTITUTO DE CIÊNCIAS MATEMÁTICAS E DE COMPUTAÇÃO}{INSTITUTE OF MATHEMATICAL AND COMPUTER SCIENCES}\par}
\vspace{0.4cm}
{\normalsize \tr%
  {Mestrado Profissional em Matemática, Estatística e Computação Aplicadas à Indústria (MECAI)}%
  {Professional Master's Program in Mathematics, Statistics and Computing Applied to Industry (MECAI)}\par}
{\normalsize \tr{Ênfase em Ciência de Dados}{Emphasis in Data Science}\par}
\vspace{3.5cm}
{\large Felipe Ramires Terrazas\par}
\vspace{2.5cm}
{\Large\bfseries\color{destaque}\titulo\par}
\vspace{1.2cm}
{\large \tr{Pré-projeto de pesquisa}{Research pre-proposal}\par}
\vspace{0.6cm}
{\normalsize \tr{Linha de pesquisa: Modelagem de Risco}{Research line: Risk Modeling}\par}
{\normalsize \tr{Orientador(a): a definir pela coordenação}{Advisor: to be assigned by the program}\par}
\vfill
{\normalsize São Carlos\par}
{\normalsize 2027\par}
\end{titlepage}

\setcounter{page}{2}
\section*{\tr{Resumo}{Abstract}}
\addcontentsline{toc}{section}{\tr{Resumo}{Abstract}}

\pt{%
Acidentes em operações de pátio logístico, como colisões de caminhões de manobra com carretas, docas e estruturas, são eventos raros, de alto custo e distribuídos de forma desigual entre veículos e unidades operacionais. Na prática, a gestão de segurança costuma comparar unidades por contagens brutas de acidentes, sem considerar quanto cada unidade operou. Como o volume de operação muda muito entre unidades e ao longo do tempo, esses rankings confundem risco com exposição e oscilam por acaso. Este projeto propõe modelar a frequência de acidentes por veículo e mês com modelos hierárquicos de contagem (Poisson e binomial negativa com efeitos aleatórios), usando como exposição as horas e distâncias registradas pela telemetria embarcada e, como preditores, taxas de eventos de condução registrados pelos sensores, como frenagens bruscas, curvas acentuadas e excesso de velocidade. O projeto também propõe estimar a subnotificação de acidentes, comparando colisões detectadas pelos sensores com os relatos formais. Os modelos serão avaliados fora da amostra, por desempenho preditivo e pela estabilidade dos rankings de risco. Caso a telemetria não possa ser usada, o mesmo arcabouço é aplicado por unidade e mês, com horas trabalhadas como exposição. Como extensão opcional, a previsão probabilística da exposição futura será combinada ao modelo de risco para projetar acidentes esperados e apoiar o planejamento da supervisão de segurança.
}
\en{%
Accidents in logistics yard operations, such as collisions of yard trucks with trailers, docks and structures, are rare, costly events that are unevenly distributed across vehicles and operating sites. In practice, safety management often compares sites by raw accident counts, ignoring how much each site operated. Since operating volume varies widely across sites and over time, such rankings confuse risk with exposure and fluctuate by chance. This project proposes to model accident frequency per vehicle and month with hierarchical count models (Poisson and negative binomial with random effects), using the hours and distances recorded by on-board telematics as exposure and, as predictors, rates of driving events recorded by the sensors, such as harsh braking, harsh turns and speeding. The project also proposes to estimate accident underreporting by comparing collisions detected by the sensors with formal reports. Models will be evaluated out of sample, by predictive performance and by the stability of risk rankings. If telematics cannot be used, the same framework is applied per site and month, with hours worked as exposure. As an optional extension, a probabilistic forecast of future exposure will be combined with the risk model to project expected accidents and support the planning of safety supervision.
}

\medskip
\noindent\textbf{\tr{Palavras-chave}{Keywords}:} \tr%
  {modelos de contagem; modelos hierárquicos; exposição ao risco; telemetria veicular; segurança do trabalho; subnotificação.}%
  {count models; hierarchical models; risk exposure; vehicle telematics; occupational safety; underreporting.}

\clearpage
\tableofcontents
\clearpage
\section{\tr{Introdução e justificativa}{Introduction and motivation}}
\label{sec:introducao}

\pt{%
Centros de distribuição e plantas industriais dependem de operações de pátio: a movimentação de carretas entre vagas, docas e portarias, feita por caminhões de manobra. É uma atividade de baixa velocidade, mas com manobras frequentes em espaço restrito, ré, engates e circulação de pedestres. Colisões com carretas, docas, portões e outros veículos geram danos materiais, interrupção da operação e, em parte dos casos, lesões.
}
\en{%
Distribution centers and industrial plants depend on yard operations: moving trailers between parking spots, docks and gates with yard trucks. The activity runs at low speed, but involves frequent maneuvers in tight spaces, reversing, coupling and pedestrian traffic. Collisions with trailers, docks, gates and other vehicles cause property damage, operational downtime and, in some cases, injuries.
}

\pt{%
O autor atua como gestor de inteligência artificial e analytics numa empresa que presta esse tipo de serviço em dezenas de unidades. Nos últimos anos, a empresa passou a ter três fontes de dados que raramente estão juntas: (i) relatos formais de acidentes, registrados num aplicativo próprio; (ii) telemetria embarcada nos caminhões, com distância, tempo de operação e eventos de condução detectados por sensores; e (iii) horas trabalhadas por unidade, vindas do sistema de folha de pagamento. O sistema de manutenção, que registra os reparos por dano em cada veículo, completa o quadro. Hoje, essas fontes são analisadas separadamente, e a comparação entre unidades usa contagens brutas de acidentes.
}
\en{%
The author works as the head of artificial intelligence and analytics at a company that provides this kind of service at dozens of sites. In recent years, the company has come to hold three data sources that are rarely available together: (i) formal accident reports, recorded in an in-house application; (ii) on-board telematics in the trucks, with distance, operating time and driving events detected by sensors; and (iii) hours worked per site, from the payroll system. The maintenance system, which records damage repairs on each vehicle, completes the picture. Today these sources are analyzed separately, and sites are compared by raw accident counts.
}

\pt{%
Contagens brutas são um indicador fraco por três motivos. Primeiro, a exposição varia muito: uma unidade com o triplo de horas operadas tende a ter mais acidentes sem ser mais arriscada. Segundo, com eventos raros, a variação por acaso é grande; uma unidade pequena com um acidente num mês pode parecer a pior da empresa \parencite{spiegelhalter2005}. Terceiro, nem todo acidente é relatado, e a propensão a relatar pode variar entre unidades, de modo que uma unidade com poucos relatos pode estar apenas relatando menos.
}
\en{%
Raw counts are a weak indicator for three reasons. First, exposure varies widely: a site with three times the operating hours tends to have more accidents without being riskier. Second, with rare events, chance variation is large; a small site with one accident in a month can look like the worst in the company \parencite{spiegelhalter2005}. Third, not every accident is reported, and the propensity to report may vary across sites, so a site with few reports may simply be reporting less.
}

\pt{%
A literatura de segurança viária trata problemas parecidos há décadas, com modelos de contagem com exposição, superdispersão e efeitos aleatórios \parencite{lord2010,mannering2014}, e a estatística aplicada oferece ferramentas para estabilizar taxas de unidades pequenas, como o bayes empírico \parencite{clayton1987,hauer1997}. Estudos com dados de condução naturalística mostram que eventos cinemáticos, como frenagens bruscas, se associam ao risco de colisão \parencite{guo2013}. Ainda assim, aplicações a operações de pátio, com veículos que andam pouco e manobram muito, e com relatos sujeitos a subnotificação, são pouco exploradas \preencher{conferir com busca bibliográfica sistemática}.
}
\en{%
Road safety research has addressed similar problems for decades, with count models that include exposure, overdispersion and random effects \parencite{lord2010,mannering2014}, and applied statistics offers tools to stabilize rates for small units, such as empirical Bayes \parencite{clayton1987,hauer1997}. Studies with naturalistic driving data show that kinematic events, such as harsh braking, are associated with crash risk \parencite{guo2013}. Still, applications to yard operations, where vehicles travel little and maneuver a lot, and where reports are subject to underreporting, remain little explored \preencher{check with a systematic literature search}.
}

\pt{%
O projeto se justifica, portanto, em duas frentes. Na prática, entrega à empresa um indicador de risco ajustado por exposição, estável e com incerteza explícita, para orientar treinamento, inspeção e supervisão. No método, combina modelos hierárquicos de contagem, exposição medida por telemetria e estimação de subnotificação num mesmo arcabouço, aplicado a um tipo de operação pouco estudado.
}
\en{%
The project is therefore justified on two fronts. In practice, it gives the company an exposure-adjusted, stable risk indicator with explicit uncertainty to guide training, inspection and supervision. Methodologically, it combines hierarchical count models, telematics-based exposure and underreporting estimation in a single framework, applied to a little-studied type of operation.
}

\section{\tr{Objetivos e hipóteses}{Objectives and hypotheses}}
\label{sec:objetivos}

\subsection{\tr{Objetivo geral}{General objective}}

\pt{%
Desenvolver e avaliar modelos estatísticos para estimar o risco de acidentes em operações de pátio, ajustado pela exposição e com incerteza quantificada, por veículo e por unidade operacional, combinando relatos de acidentes, telemetria veicular e horas trabalhadas.
}
\en{%
To develop and evaluate statistical models that estimate accident risk in yard operations, adjusted for exposure and with quantified uncertainty, per vehicle and per operating site, by combining accident reports, vehicle telematics and hours worked.
}

\subsection{\tr{Objetivos específicos}{Specific objectives}}

\begin{enumerate}[label=(\alph*)]
  \item \tr%
    {Integrar as três fontes num painel veículo $\times$ mês, com regras de ligação documentadas e qualidade da ligação medida em amostra revisada manualmente.}%
    {Integrate the three sources into a vehicle $\times$ month panel, with documented linkage rules and linkage quality measured on a manually reviewed sample.}
  \item \tr%
    {Comparar modelos de contagem (Poisson, binomial negativa, com e sem inflação de zeros, com e sem efeitos aleatórios de veículo e unidade) quanto ao ajuste e ao desempenho preditivo fora da amostra.}%
    {Compare count models (Poisson, negative binomial, with and without zero inflation, with and without vehicle and site random effects) in terms of fit and out-of-sample predictive performance.}
  \item \tr%
    {Medir quanto as taxas de eventos de condução registradas pela telemetria acrescentam à previsão de acidentes, além da exposição e da unidade.}%
    {Measure how much the driving event rates recorded by telematics add to accident prediction beyond exposure and site.}
  \item \tr%
    {Estimar a probabilidade de que uma colisão detectada pelos sensores seja relatada formalmente, e sua variação entre unidades.}%
    {Estimate the probability that a sensor-detected collision is formally reported, and how it varies across sites.}
  \item \tr%
    {Produzir rankings de risco por unidade ajustados por exposição, com intervalos de incerteza, e compará-los aos rankings por contagem bruta quanto à estabilidade ao longo do tempo.}%
    {Produce exposure-adjusted site risk rankings with uncertainty intervals, and compare them with raw-count rankings in terms of stability over time.}
  \item \textit{\tr{(Opcional, extensão)}{(Optional, extension)}} \tr%
    {Prever de forma probabilística a exposição (horas de operação) de cada unidade nos meses seguintes, como insumo para o planejamento de escala e de supervisão, e combinar essa previsão ao modelo de risco para projetar acidentes esperados em níveis agregados (região e trimestre).}%
    {Forecast each site's exposure (operating hours) for the following months probabilistically, as input for staffing and supervision planning, and combine this forecast with the risk model to project expected accidents at aggregate levels (region and quarter).}
\end{enumerate}

\pt{%
O objetivo (f) é uma extensão: será desenvolvido se os objetivos (a) a (e) estiverem concluídos dentro do cronograma (Seção~\ref{sec:cronograma}).
}
\en{%
Objective (f) is an extension: it will be pursued if objectives (a) to (e) are completed on schedule (Section~\ref{sec:cronograma}).
}

\subsection{\tr{Hipóteses}{Hypotheses}}

\begin{description}[leftmargin=1.2cm,style=nextline]
  \item[H1] \tr%
    {Taxas por unidade estimadas por modelos hierárquicos com exposição preveem a taxa de acidentes do período seguinte melhor que taxas brutas, e, ao contrário das contagens brutas, não se confundem com o tamanho da unidade.}%
    {Site rates estimated by hierarchical models with exposure predict the next period's accident rate better than raw rates and, unlike raw counts, are not confounded with site size.}
  \item[H2] \tr%
    {As taxas de eventos de condução melhoram o desempenho preditivo fora da amostra (escore logarítmico) em relação a um modelo apenas com exposição e unidade.}%
    {Driving event rates improve out-of-sample predictive performance (log score) relative to a model with exposure and site only.}
  \item[H3] \tr%
    {A probabilidade de relato de colisões detectadas é menor que um e varia entre unidades, de modo que as contagens de relatos subestimam o risco de forma desigual.}%
    {The reporting probability of detected collisions is below one and varies across sites, so that report counts underestimate risk unevenly.}
\end{description}

\section{\tr{Fundamentação teórica}{Background}}
\label{sec:fundamentacao}

\subsection{\tr{Modelos de contagem com exposição}{Count models with exposure}}

\pt{%
Seja $Y_{it}$ o número de acidentes do veículo $i$ no mês $t$ e $E_{it}$ sua exposição (horas de operação ou distância). O modelo de Poisson com \textit{offset} supõe $Y_{it} \sim \text{Poisson}(\mu_{it})$, com $\log \mu_{it} = \log E_{it} + \mathbf{x}_{it}^\top \bm{\beta}$, de modo que $\exp(\mathbf{x}_{it}^\top \bm{\beta})$ é uma taxa por unidade de exposição \parencite{cameron2013}. Dados de acidentes costumam ter variância maior que a média (superdispersão), o que leva à binomial negativa, e excesso de zeros, que pode pedir modelos com inflação de zeros \parencite{lambert1992}. A escolha entre essas alternativas em dados de acidentes é discutida em detalhe por \textcite{lord2010}; em segurança viária, a inflação de zeros é criticada por supor um estado de risco nulo sem interpretação física, quando o excesso de zeros costuma vir de exposição pequena ou de heterogeneidade não modelada \parencite{lord2005}.
}
\en{%
Let $Y_{it}$ be the number of accidents of vehicle $i$ in month $t$ and $E_{it}$ its exposure (operating hours or distance). The Poisson model with an offset assumes $Y_{it} \sim \text{Poisson}(\mu_{it})$, with $\log \mu_{it} = \log E_{it} + \mathbf{x}_{it}^\top \bm{\beta}$, so that $\exp(\mathbf{x}_{it}^\top \bm{\beta})$ is a rate per unit of exposure \parencite{cameron2013}. Accident data usually show variance larger than the mean (overdispersion), which leads to the negative binomial, and excess zeros, which may call for zero-inflated models \parencite{lambert1992}. The choice among these alternatives for crash data is discussed in detail by \textcite{lord2010}; in road safety, zero inflation is criticized for assuming a zero-risk state with no physical interpretation, when excess zeros usually come from small exposure or unmodeled heterogeneity \parencite{lord2005}.
}

\subsection{\tr{Estrutura hierárquica e encolhimento}{Hierarchical structure and shrinkage}}

\pt{%
Veículos estão agrupados em unidades, e observações do mesmo veículo se repetem no tempo. Efeitos aleatórios de veículo e de unidade capturam essa dependência e produzem estimativas que encolhem as taxas de grupos pequenos em direção à média geral, na medida da informação disponível \parencite{gelman2007}. É a mesma ideia do Bayes empírico usado em mapeamento de doenças e em segurança viária para identificar locais de alto risco sem ser enganado pela regressão à média \parencite{clayton1987,hauer1997}. Rankings de instituições têm limites conhecidos, porque a incerteza sobre a posição de cada unidade costuma ser grande \parencite{goldstein1996}, e gráficos de funil oferecem uma forma visual de comparar unidades com tamanhos diferentes \parencite{spiegelhalter2005}. Os modelos podem ser ajustados por máxima verossimilhança com aproximação de Laplace \parencite{brooks2017} ou de forma bayesiana \parencite{burkner2017}.
}
\en{%
Vehicles are grouped in sites, and observations of the same vehicle repeat over time. Vehicle and site random effects capture this dependence and yield estimates that shrink the rates of small groups toward the overall mean, to the extent of the available information \parencite{gelman2007}. This is the same idea as the empirical Bayes approach used in disease mapping and in road safety to identify high-risk locations without being misled by regression to the mean \parencite{clayton1987,hauer1997}. Institutional rankings have known limits, because the uncertainty about each unit's position is usually large \parencite{goldstein1996}, and funnel plots offer a visual way to compare units of different sizes \parencite{spiegelhalter2005}. The models can be fitted by maximum likelihood with the Laplace approximation \parencite{brooks2017} or in a Bayesian way \parencite{burkner2017}.
}

\subsection{\tr{Telemetria e risco}{Telematics and risk}}

\pt{%
Sistemas de telemetria registram eventos cinemáticos (acelerações, frenagens e curvas acima de limiares) e eventos de comportamento (velocidade acima do limite, uso de cinto, uso de celular). Estudos de condução naturalística associam a frequência desses eventos ao risco de colisão de cada motorista \parencite{guo2013}. Em caminhões, eventos críticos registrados pela telemetria de uma grande transportadora dos EUA se associaram a acidentes num modelo binomial negativo com a distância como exposição \parencite{cai2021}, e frenagens bruscas por hora previram sinistros com dano material \parencite{driessen2024}. Em pátio, onde as velocidades são baixas e as manobras dominam, não é óbvio quais eventos carregam sinal; o projeto testa isso diretamente (H2).
}
\en{%
Telematics systems record kinematic events (accelerations, braking and turns above thresholds) and behavioral events (speeding, seat belt use, phone use). Naturalistic driving studies associate the frequency of these events with each driver's crash risk \parencite{guo2013}. For trucks, safety-critical events recorded by the telematics of a large US carrier were associated with crashes in a negative binomial model with distance as exposure \parencite{cai2021}, and harsh braking per hour predicted property damage incidents \parencite{driessen2024}. In a yard, where speeds are low and maneuvering dominates, it is not obvious which events carry signal; the project tests this directly (H2).
}

\subsection{\tr{Subnotificação}{Underreporting}}

\pt{%
A subnotificação de lesões de trabalho é grande e varia entre locais: estudos nos EUA que cruzam o levantamento oficial com outros registros encontraram grande parte das lesões fora das estatísticas, com muita variação entre estados \parencite{leigh2004,rosenman2006,boden2008}. Quando o mesmo tipo de evento é registrado por duas fontes imperfeitas, a sobreposição entre elas informa quantos eventos nenhuma das duas captou. É o princípio da captura e recaptura com múltiplas fontes, usado em epidemiologia \parencite{iwgdmf1995}. Aqui, as fontes são os relatos formais e as ordens de manutenção abertas por dano. A hipótese de independência entre as fontes raramente vale na prática: com dependência positiva, como a de acidentes graves que tendem a ser ao mesmo tempo relatados e reparados, o total é subestimado \parencite{hook1995}, e o projeto tratará isso de forma explícita (Seção~\ref{sec:subnotificacao}). A ligação entre registros das duas fontes segue a lógica da ligação probabilística de registros \parencite{fellegi1969}.
}
\en{%
Underreporting of work injuries is large and varies across places: US studies cross-matching the official survey with other records found a large share of injuries missing from the statistics, with much variation across states \parencite{leigh2004,rosenman2006,boden2008}. When the same type of event is recorded by two imperfect sources, the overlap between them informs how many events neither source captured. This is the principle of multiple-source capture-recapture used in epidemiology \parencite{iwgdmf1995}. Here the sources are formal reports and maintenance orders opened for damage. The independence assumption between sources rarely holds in practice: with positive dependence, such as severe accidents tending to be both reported and repaired, the total is underestimated \parencite{hook1995}, and the project will handle this explicitly (Section~\ref{sec:subnotificacao}). Linking records across the two sources follows the logic of probabilistic record linkage \parencite{fellegi1969}.
}

\subsection{\tr{Previsão probabilística}{Probabilistic forecasting}}

\pt{%
Para a extensão opcional, a exposição futura será prevista como distribuição, não como valor pontual, com métodos de séries temporais \parencite{hyndman2021} e avaliada por regras de pontuação próprias, como o CRPS e a perda quantílica \parencite{gneiting2007}.
}
\en{%
For the optional extension, future exposure will be forecast as a distribution rather than a point value, with time series methods \parencite{hyndman2021}, and evaluated with proper scoring rules such as the CRPS and the quantile loss \parencite{gneiting2007}.
}

\section{\tr{Dados}{Data}}
\label{sec:dados}

\pt{%
O estudo usa apenas dados que a empresa já coleta na rotina; não há coleta própria. Os números abaixo são ordens de grandeza, de um levantamento preliminar feito pelo autor, e serão atualizados no início do projeto.
}
\en{%
The study uses only data the company already collects routinely; there is no new data collection. The figures below are orders of magnitude from a preliminary survey by the author and will be updated at the start of the project.
}

\begin{table}[h]
\centering
\small
\caption{\tr{Fontes de dados previstas}{Planned data sources}}
\label{tab:fontes}
\begin{tabularx}{\textwidth}{>{\raggedright\arraybackslash}p{3.2cm} >{\raggedright\arraybackslash}X >{\raggedright\arraybackslash}p{3.6cm}}
\toprule
\tr{Fonte}{Source} & \tr{Conteúdo}{Content} & \tr{Escala aproximada}{Approximate scale} \\
\midrule
\tr{Relatos de acidentes}{Accident reports} &
\tr{Data, tipo de evento, unidade, número do veículo, gravidade, descrição}%
   {Date, event type, site, vehicle number, severity, description} &
\tr{Centenas de relatos em pouco mais de um ano}{Hundreds of reports over a little more than a year} \\
\tr{Telemetria veicular}{Vehicle telematics} &
\tr{Viagens (distância, duração), eventos de condução, colisões detectadas, unidade do veículo}%
   {Trips (distance, duration), driving events, detected collisions, vehicle site} &
\tr{Cerca de 300 veículos; mais de um milhão de viagens}{About 300 vehicles; over one million trips} \\
\tr{Horas trabalhadas}{Hours worked} &
\tr{Horas por unidade e mês}{Hours per site and month} &
\tr{Dezenas de unidades}{Dozens of sites} \\
\tr{Manutenção}{Maintenance} &
\tr{Ordens de reparo por veículo, com motivo (dano, acidente, preventiva)}%
   {Repair orders per vehicle, with reason (damage, accident, preventive)} &
\tr{Centenas de ordens por dano}{Hundreds of damage orders} \\
\bottomrule
\end{tabularx}
\end{table}

\pt{%
Um levantamento preliminar indicou que a ligação entre as fontes é viável: o número do veículo informado nos relatos coincide com o identificador da telemetria na maior parte dos casos, e o cadastro da telemetria associa cada veículo a uma unidade com o mesmo código usado no sistema de folha. O período comum às três fontes começa em 2025 \preencher{atualizar período e cobertura no início do projeto}.
}
\en{%
A preliminary survey indicated that linking the sources is feasible: the vehicle number given in the reports matches the telematics identifier in most cases, and the telematics registry assigns each vehicle to a site using the same code as the payroll system. The period common to the three sources starts in 2025 \preencher{update period and coverage at the start of the project}.
}

\pt{%
Os problemas de qualidade já conhecidos fazem parte do trabalho: relatos de teste misturados aos reais no período de implantação do aplicativo, unidade e veículo digitados em texto livre em parte dos relatos, viagens com registro de fim inválido na telemetria e mudanças de veículos entre unidades ao longo do tempo. Cada regra de limpeza será documentada e seu efeito sobre os resultados, testado em análise de sensibilidade.
}
\en{%
Known data quality issues are part of the work: test reports mixed with real ones during the application roll-out, site and vehicle typed as free text in some reports, trips with invalid end records in the telematics, and vehicles moving between sites over time. Every cleaning rule will be documented and its effect on the results tested in a sensitivity analysis.
}

\subsection*{\tr{Descrição preliminar}{Preliminary description}}

\pt{%
Uma análise descritiva agregada dos dados da operação, só com valores relativos (Figura~\ref{fig:desc-pp}), já mostra quatro fatos que orientam o desenho. (a)~As horas de viagem por veículo-mês são parecidas nos três tipos de uso, mas as milhas, não: o spotter anda a cerca de \descVelSpotter{} da velocidade do veículo local, o que justifica o tempo como exposição. (b)~Os acidentes parecem concentrados em poucos veículos (Gini de \descGiniObs), mas quase toda essa concentração é a esperada por acaso e exposição (\descGiniPoisson{} sob Poisson); concentração não é, por si só, evidência de risco maior. (c)~A hora do acidente está registrada em \descComHora\% dos relatos, e os acidentes se distribuem pelos três turnos, o que torna possíveis análises por turno. (d)~Relatos, ordens de manutenção por dano e colisões detectadas variam de forma diferente ao longo do ano, sinal de que captam universos diferentes. A dissertação terá um capítulo descritivo completo, com a distribuição de cada métrica usada nos modelos, seguido dos capítulos preditivo e prescritivo (Seção~\ref{sec:metodologia}).
}
\en{%
An aggregated descriptive analysis of the operational data, with relative values only (Figure~\ref{fig:desc-pp}), already shows four facts that shape the design. (a)~Trip hours per vehicle-month are similar across the three use types, but miles are not: the spotter moves at about \descVelSpotter{} of the local vehicle's speed, which justifies time as exposure. (b)~Accidents look concentrated in a few vehicles (Gini of \descGiniObs), but almost all of this concentration is what chance and exposure produce (\descGiniPoisson{} under Poisson); concentration is not, in itself, evidence of higher risk. (c)~The accident time is recorded in \descComHora\% of reports, and accidents are spread over the three shifts, which makes shift-level analyses possible. (d)~Reports, damage repair orders and detected collisions vary differently over the year, a sign that they capture different universes. The dissertation will have a full descriptive chapter, with the distribution of each metric used in the models, followed by the predictive and prescriptive chapters (Section~\ref{sec:metodologia}).
}

\begin{figure}[!htbp]
\centering
\includegraphics[width=\textwidth]{\arqfig{fig_desc_pp}}
\caption{\tr%
  {Descrição preliminar dos dados da operação (só valores relativos). (a)~Horas de viagem por veículo-mês, por tipo de uso, relativas à mediana. (b)~Fração acumulada dos acidentes e das horas contra a fração acumulada de veículos, ordenados do maior para o menor; faixa: 95\% das curvas simuladas sob Poisson com taxa constante por tipo de uso. (c)~Distribuição dos acidentes por hora do dia; faixas: turnos de 8 horas. (d)~Relatos de acidente, ordens de manutenção por dano e colisões detectadas pela telemetria, por mês, relativos à média de cada série.}%
  {Preliminary description of the operational data (relative values only). (a)~Trip hours per vehicle-month, by use type, relative to the median. (b)~Cumulative share of accidents and hours against the cumulative share of vehicles, ordered from largest to smallest; band: 95\% of curves simulated under Poisson with a constant rate per use type. (c)~Distribution of accidents by hour of day; bands: 8-hour shifts. (d)~Accident reports, damage repair orders and telematics-detected collisions per month, relative to each series' mean.}}
\label{fig:desc-pp}
\end{figure}

\begin{nota}[\tr{Privacidade}{Privacy}]
\tr%
  {Nenhum dado identifica pessoas. A análise usa veículo, unidade e mês; identificadores de motorista, nomes, imagens de câmera e trajetos de GPS não são usados. Os dados não saem do ambiente da empresa; apenas tabelas agregadas entram no repositório da pesquisa e na dissertação, que não nomeará a empresa, seus clientes nem suas unidades.}%
  {No data identifies people. The analysis uses vehicle, site and month; driver identifiers, names, camera footage and GPS tracks are not used. Data stays within the company's environment; only aggregated tables enter the research repository and the dissertation, which will not name the company, its clients or its sites.}
\end{nota}

\section{\tr{Metodologia}{Methodology}}
\label{sec:metodologia}

\subsection{\tr{Etapa 1: integração e painel}{Stage 1: integration and panel}}

\pt{%
As três fontes serão integradas num painel com uma linha por veículo e mês. Para cada linha: o número de acidentes relatados, a exposição (horas com motor ligado e distância percorrida), as contagens de cada tipo de evento de condução, características do veículo (tipo de uso, combustível) e a unidade em que o veículo estava. A ligação entre relatos e veículos usa regras determinísticas (normalização do número do veículo) e, para os casos restantes, comparação aproximada inspirada em \textcite{fellegi1969}. A qualidade será medida numa amostra revisada manualmente (precisão e cobertura da ligação), e os modelos serão reajustados só com ligações exatas, como teste de sensibilidade.
}
\en{%
The three sources will be integrated into a panel with one row per vehicle and month. For each row: the number of reported accidents, exposure (engine hours and distance traveled), the counts of each type of driving event, vehicle characteristics (use type, fuel) and the site where the vehicle was. Linking reports to vehicles uses deterministic rules (normalizing the vehicle number) and, for the remaining cases, approximate matching inspired by \textcite{fellegi1969}. Linkage quality will be measured on a manually reviewed sample (precision and coverage), and models will be refitted with exact links only as a sensitivity check.
}

\subsection{\tr{Etapa 2: modelos de risco}{Stage 2: risk models}}

\pt{%
O modelo principal é uma binomial negativa hierárquica:
}
\en{%
The main model is a hierarchical negative binomial:
}
\begin{equation}
\label{eq:modelo}
\begin{aligned}
Y_{it} &\sim \text{NB}(\mu_{it}, \phi), \\
\log \mu_{it} &= \log E_{it} + \mathbf{x}_{it}^\top \bm{\beta} + \gamma_{t} + u_{s(i,t)} + v_i, \\
u_s &\sim \mathcal{N}(0, \sigma_u^2), \qquad v_i \sim \mathcal{N}(0, \sigma_v^2),
\end{aligned}
\end{equation}
\pt{%
em que $s(i,t)$ é a unidade do veículo $i$ no mês $t$, $\mathbf{x}_{it}$ reúne as taxas de eventos de condução por unidade de exposição e as características do veículo, $\gamma_t$ é um efeito de calendário (tendência e sazonalidade) e $\phi$ é o parâmetro de dispersão. O modelo será comparado a alternativas mais simples (Poisson, sem efeitos aleatórios, sem telemetria) e a versões com inflação de zeros. O diagnóstico usará resíduos quantílicos aleatorizados \parencite{dunn1996}.
}
\en{%
where $s(i,t)$ is the site of vehicle $i$ in month $t$, $\mathbf{x}_{it}$ holds the driving event rates per unit of exposure and the vehicle characteristics, $\gamma_t$ is a calendar effect (trend and seasonality) and $\phi$ is the dispersion parameter. The model will be compared with simpler alternatives (Poisson, without random effects, without telematics) and with zero-inflated versions. Diagnostics will use randomized quantile residuals \parencite{dunn1996}.
}

\pt{%
Os coeficientes de telemetria serão lidos como associações, não como efeitos causais: o mesmo motorista, a mesma unidade e o mesmo tipo de manobra podem gerar ao mesmo tempo mais eventos e mais acidentes. A definição de \enquote{acidente} também será fixada antes da modelagem (por exemplo, separar colisões com dano material de lesões e de quase-acidentes) e testada com definições alternativas.
}
\en{%
Telematics coefficients will be read as associations, not causal effects: the same driver, the same site and the same type of maneuver can generate both more events and more accidents. The definition of \enquote{accident} will also be fixed before modeling (for example, separating property-damage collisions from injuries and near misses) and tested with alternative definitions.
}

\subsection{\tr{Etapa 3: validação}{Stage 3: validation}}

\pt{%
A avaliação será fora da amostra, de duas formas: validação temporal (ajuste nos meses iniciais, previsão dos últimos meses) e validação por unidade (deixando uma unidade de fora por vez). As métricas serão o escore logarítmico e a calibração da distribuição prevista \parencite{gneiting2007}, além da fração de acidentes concentrada nos veículos e unidades classificados como de maior risco. Para H1, cada indicador por unidade (contagem bruta, taxa bruta e taxa estimada pelo modelo~\eqref{eq:modelo}) calculado num semestre será usado para prever a taxa observada no semestre seguinte, com erro ponderado pela exposição; a correlação de cada indicador com o tamanho da unidade também será reportada, porque contagens brutas são estáveis justamente por refletirem o tamanho.
}
\en{%
Evaluation will be out of sample, in two ways: temporal validation (fit on the early months, predict the last months) and site-level validation (leaving one site out at a time). Metrics will be the log score and the calibration of the predictive distribution \parencite{gneiting2007}, plus the share of accidents concentrated in the vehicles and sites ranked as highest risk. For H1, each site indicator (raw count, raw rate and rate estimated by model~\eqref{eq:modelo}) computed in one semester will be used to predict the rate observed in the next semester, with exposure-weighted error; the correlation of each indicator with site size will also be reported, since raw counts are stable precisely because they reflect size.
}

\subsection{\tr{Etapa 4: subnotificação}{Stage 4: underreporting}}
\label{sec:subnotificacao}

\pt{%
Cada colisão detectada pelos sensores será comparada aos relatos do mesmo veículo numa janela de tempo curta. Com isso, estima-se a probabilidade de relato $p_s$ de uma colisão detectada, por unidade, com um modelo logístico hierárquico. Dois cuidados são centrais: os sensores não detectam todas as colisões (impactos leves em baixa velocidade podem não disparar o alarme), e parte das detecções são falsos alarmes. Por isso, a análise será restrita aos eventos revisados como colisão no próprio sistema de telemetria, e a estimativa será apresentada como limite para a subnotificação, não como valor exato. Se houver dados suficientes, um modelo de captura e recaptura com duas fontes e dependência entre elas \parencite{iwgdmf1995} será usado para estimar o total de colisões.
}
\en{%
Each sensor-detected collision will be compared with reports for the same vehicle within a short time window. This yields an estimate of the reporting probability $p_s$ of a detected collision, per site, with a hierarchical logistic model. Two caveats are central: sensors do not detect every collision (light, low-speed impacts may not trigger the alarm), and some detections are false alarms. The analysis will therefore be restricted to events reviewed as collisions in the telematics system itself, and the estimate will be presented as a bound on underreporting, not an exact value. If data allow, a two-source capture-recapture model with dependence between sources \parencite{iwgdmf1995} will be used to estimate the total number of collisions.
}

\subsection{\tr{Plano alternativo}{Fallback plan}}

\pt{%
Se o uso da telemetria for restrito, o mesmo arcabouço é aplicado ao nível de unidade e mês, com as horas trabalhadas da folha como exposição. Os objetivos (b) e (e) e a hipótese H1 seguem testáveis; o objetivo (c) e a hipótese H2 deixam de existir, e o objetivo (d) é substituído pela comparação entre unidades com práticas de relato diferentes \preencher{detalhar se o plano alternativo for acionado}.
}
\en{%
If the use of telematics is restricted, the same framework is applied at the site-month level, with payroll hours worked as exposure. Objectives (b) and (e) and hypothesis H1 remain testable; objective (c) and hypothesis H2 are dropped, and objective (d) is replaced by a comparison between sites with different reporting practices \preencher{detail if the fallback plan is triggered}.
}

\subsection{\tr{Extensão opcional: projeção de acidentes esperados}{Optional extension: projecting expected accidents}}

\pt{%
A exposição de cada unidade nos meses seguintes será prevista como distribuição, com modelos de séries temporais hierárquicos \parencite{hyndman2021}, e avaliada por CRPS e perda quantílica \parencite{gneiting2007}. Combinando essa distribuição com o modelo~\eqref{eq:modelo} por simulação, obtém-se a distribuição de acidentes esperados por unidade e mês. O resultado é uma ferramenta para decidir onde concentrar inspeções e supervisão no mês seguinte, considerando tanto o risco por hora quanto o volume previsto.
}
\en{%
Each site's exposure for the following months will be forecast as a distribution, with hierarchical time series models \parencite{hyndman2021}, and evaluated by CRPS and quantile loss \parencite{gneiting2007}. Combining this distribution with model~\eqref{eq:modelo} by simulation yields the distribution of expected accidents per site and month. The result is a tool to decide where to focus inspections and supervision in the following month, accounting for both the risk per hour and the forecast volume.
}

\subsection{\tr{Ferramentas e reprodutibilidade}{Tools and reproducibility}}

\pt{%
A extração e a agregação rodam dentro do ambiente de dados da empresa, em SQL. A modelagem usa Python e R, com código versionado, sementes fixas e um esqueleto de análise já testado com dados sintéticos que imitam a estrutura do painel. Toda figura e tabela da dissertação será gerada por script.
}
\en{%
Extraction and aggregation run inside the company's data environment, in SQL. Modeling uses Python and R, with versioned code, fixed seeds and an analysis skeleton already tested on synthetic data that mimic the panel structure. Every figure and table in the dissertation will be generated by script.
}

\section{\tr{Resultados esperados e contribuições}{Expected results and contributions}}
\label{sec:resultados}

\begin{itemize}
  \item \tr%
    {\textbf{Método:} uma comparação sistemática de modelos de contagem com exposição para acidentes em operações de pátio, incluindo a escolha da medida de exposição e o valor preditivo da telemetria, desenhada para a estrada, num contexto de baixa velocidade e muitas manobras.}%
    {\textbf{Method:} a systematic comparison of count models with exposure for accidents in yard operations, including the choice of exposure measure and the predictive value of road-oriented telematics in a low-speed, maneuver-heavy setting.}
  \item \tr%
    {\textbf{Subnotificação:} uma estimativa, com limites explícitos, de quanto as contagens de relatos subestimam os acidentes e de como isso varia entre unidades, além de uma medida da cobertura da detecção automática de colisão no pátio.}%
    {\textbf{Underreporting:} an estimate, with explicit bounds, of how much report counts underestimate accidents and how this varies across sites, plus a measure of the coverage of automatic crash detection in the yard.}
  \item \tr%
    {\textbf{Aplicação:} um indicador de risco por unidade e veículo, ajustado por exposição e com incerteza, pronto para ser usado pela gestão de segurança, e, se a extensão for feita, uma previsão probabilística de horas de operação para planejamento de escala e supervisão.}%
    {\textbf{Application:} a site and vehicle risk indicator, adjusted for exposure and with uncertainty, ready for use by safety management, and, if the extension is carried out, a probabilistic forecast of operating hours for staffing and supervision planning.}
  \item \tr%
    {\textbf{Reprodutibilidade:} código aberto do pipeline de análise, em Python e R, já testado com dados sintéticos, com um conjunto público e com agregados dos dados da operação (Seção~\ref{sec:estudo-preliminar}), que pode ser reaplicado por outras operações com dados semelhantes.}%
    {\textbf{Reproducibility:} open-source analysis pipeline code, in Python and R, already tested on synthetic data, on a public dataset and on aggregates of the operational data (Section~\ref{sec:estudo-preliminar}), that can be reapplied by other operations with similar data.}
\end{itemize}

\pt{%
Os resultados podem ser negativos de forma útil. A exploração preliminar sugere que os eventos de condução pouco acrescentam quando se comparam veículos do mesmo tipo de uso. Se o teste de equivalência de H2 confirmar efeito desprezível, isso indica quais indicadores não vale a pena acompanhar e mostra o limite de ferramentas desenhadas para a estrada quando aplicadas ao pátio.
}
\en{%
Results may be usefully negative. The preliminary exploration suggests that driving events add little when comparing vehicles of the same use type. If the H2 equivalence test confirms a negligible effect, this shows which indicators are not worth tracking and the limits of road-oriented tools when applied to the yard.
}

\section{\tr{Viabilidade e aspectos éticos}{Feasibility and ethical aspects}}
\label{sec:viabilidade}

\pt{%
O projeto é compatível com a atividade profissional do autor, que é responsável pela área de dados da empresa e já trabalha com as fontes descritas. Os dados existem, cobrem mais de um ano e não dependem de coleta nova. O uso dos dados na pesquisa será formalizado com a empresa por autorização escrita após a seleção, com as restrições descritas na Seção~\ref{sec:dados}: análise dentro do ambiente da empresa, nenhum dado pessoal e apenas agregados na dissertação.
}
\en{%
The project is compatible with the author's professional activity, as the author leads the company's data area and already works with the sources described. The data exist, cover more than a year and do not depend on new collection. Use of the data in the research will be formalized with the company through a written authorization after selection, with the restrictions described in Section~\ref{sec:dados}: analysis inside the company's environment, no personal data and only aggregates in the dissertation.
}

\pt{%
Os principais riscos e as respostas previstas são: (i) restrição ao uso da telemetria, coberta pelo plano alternativo por unidade e mês; (ii) poucos acidentes por veículo, que justificam os modelos hierárquicos e o encolhimento; (iii) ruído nos relatos do período de implantação, tratado com regras documentadas e análise de sensibilidade. A necessidade de submissão a comitê de ética para uso de dados secundários agregados será verificada com a coordenação \preencher{confirmar exigência do programa}.
}
\en{%
The main risks and planned responses are: (i) restricted use of telematics, covered by the site-month fallback plan; (ii) few accidents per vehicle, which justify the hierarchical models and shrinkage; (iii) noise in reports from the roll-out period, handled with documented rules and sensitivity analysis. Whether submission to an ethics committee is required for aggregated secondary data will be checked with the program \preencher{confirm program requirement}.
}

\section{\tr{Cronograma}{Timeline}}
\label{sec:cronograma}

\pt{%
O prazo de depósito do programa é de 36 meses, e o cronograma está dividido em seis semestres. As disciplinas se concentram nos três primeiros, e a proficiência em inglês será comprovada nos primeiros 18 meses. Como o programa não exige exame de qualificação, a validação dos modelos (Etapa 3) funciona como ponto de controle interno ao fim do segundo ano.
}
\en{%
The program's submission deadline is 36 months, and the timeline is split into six semesters. Coursework is concentrated in the first three, and English proficiency will be certified within the first 18 months. Since the program does not require a qualifying exam, model validation (Stage 3) serves as an internal checkpoint at the end of the second year.
}

\newcommand{\marca}{$\bullet$}
\begin{table}[h]
\centering
\small
\caption{\tr{Cronograma por semestre}{Timeline by semester}}
\label{tab:cronograma}
\begin{tabularx}{\textwidth}{>{\raggedright\arraybackslash}X *{6}{c}}
\toprule
\tr{Atividade}{Activity} & S1 & S2 & S3 & S4 & S5 & S6 \\
\midrule
\tr{Disciplinas}{Coursework} & \marca & \marca & \marca & & & \\
\tr{Proficiência em inglês}{English proficiency} & \marca & \marca & \marca & & & \\
\tr{Revisão bibliográfica}{Literature review} & \marca & \marca & & & & \\
\tr{Autorização e acesso aos dados}{Authorization and data access} & \marca & & & & & \\
\tr{Etapa 1: integração e painel}{Stage 1: integration and panel} & & \marca & & & & \\
\tr{Etapa 2: modelos de risco}{Stage 2: risk models} & & \marca & \marca & \marca & & \\
\tr{Etapa 3: validação}{Stage 3: validation} & & & \marca & \marca & & \\
\tr{Etapa 4: subnotificação}{Stage 4: underreporting} & & & & \marca & \marca & \\
\tr{Extensão opcional (objetivo f)}{Optional extension (objective f)} & & & & & \marca & \\
\tr{Redação da dissertação}{Dissertation writing} & & & & \marca & \marca & \marca \\
\tr{Depósito e defesa}{Submission and defense} & & & & & & \marca \\
\bottomrule
\end{tabularx}
\end{table}


\begingroup
\singlespacing
\setlength{\bibitemsep}{\baselineskip}
\printbibliography[title={\tr{Referências}{References}}]
\endgroup

\end{document}
